% §1 Introduction

\section{Introduction}
\label{sec:intro}

Large language models (LLMs) vary widely in capability and inference cost.
Deploying a single model for every query is inefficient: some inputs are answered correctly by a small model, while others require a larger one \citep{hu2024routerbench,ding2024hybrid,chen2023frugalgpt}.
\textbf{Multi-LLM routing} assigns each query to a model in a fixed pool $\mathcal{M}$ so as to trade off response quality against inference cost \citep{ong2024routellm}.

\paragraph{Problem.}
The routing objective is \textbf{appropriate model selection}: route to the \textbf{weakest model expected to answer correctly}, escalating only when additional capability is likely to help.
Most prior work develops \textbf{routing algorithms}---supervised routers \citep{ong2024routellm,ding2024hybrid,feng2025graphrouter}, learned pre-hoc estimators \citep{cao2026scope}, or post-generation selectors \citep{guha2024smoothie,niu2026cascal}.
We instead ask a prior question: \textbf{what routing-relevant information is available before any routing decision is made?}
Specifically, can \textbf{unsupervised pre-inference signals}---statistics derived from query text and prefill probes before full generation---support appropriate model selection?

\paragraph{Research question.}
\textit{Can unsupervised pre-inference signals support selecting the appropriate LLM from a fixed pool before generation?}

We do not ask which isolated statistic achieves the highest AUROC.
We ask which \textbf{latent routing dimensions}---task difficulty, model uncertainty, model disagreement, escalation potential---are measurable before generation, and whether a calibrated policy can use them (Table~\ref{tab:dimensions}).

\paragraph{Approach.}
We treat routing need as \textbf{latent}: each dimension is \textbf{operationalized} by one prefill statistic (e.g., task difficulty by \texttt{piece\_count}; escalation potential by $\Delta m_{\mathrm{gain}}$).
Studies I--III ask \textbf{which dimension predicts routing opportunity} and whether dimensions encode \textbf{different} aspects of that need.
Signal characterization supplies \textbf{evidence}; routing evaluation supplies \textbf{validation} (\S\ref{sec:routing}).

\paragraph{What ``unsupervised'' means.}
\textbf{Signal extraction} is unsupervised at inference time.
Offline oracle labels evaluate whether each dimension predicts opportunity (Studies I--III).
Routing evaluation fits a demonstration policy on CALIB only (\S\ref{sec:method:supervision}).

\paragraph{Contributions.}
\begin{enumerate}
  \item An empirical study of \textbf{unsupervised pre-inference signals} for multi-LLM routing, with appropriate selection defined over a fixed pool and no routing labels at extraction time.
  \item A dimensional account of routing need: four \textbf{latent routing dimensions} with prefill operationalizations (Table~\ref{tab:dimensions}).
  \item \textbf{Empirical evidence} on ARC-Challenge: multiple dimensions predict opportunity; \textbf{task difficulty and model uncertainty} overlap, while \textbf{model disagreement and escalation potential} separate routable from irrecoverable hard queries; dimensions are partially complementary (\S\ref{sec:characterization}).
  \item \textbf{Routing evaluation}: a calibrated logistic policy leaves exploitable headroom relative to an offline oracle (\S\ref{sec:routing}).
\end{enumerate}

\paragraph{Main findings (ARC test, $n{=}1{,}172$).}
The answer is \textbf{partially}.
\begin{enumerate}
  \item \textbf{Dimensions predict opportunity.} Routing opportunity occurs on 43.3\% of queries; escalation potential reaches AUROC ${\sim}0.61$ (Table~\ref{tab:study2}).
  \item \textbf{Dimensions differ.} Model uncertainty does not separate opportunity from too-hard queries ($d{\approx}0.03$); escalation potential does ($d{\approx}0.72$).
  \item \textbf{Routing leaves headroom.} A CALIB-fit calibrated policy matches always-strong (69.2\%); oracle 74.4\% (Table~\ref{tab:routing}).
\end{enumerate}

\paragraph{Paper organization.}
\S\ref{sec:related}--\S\ref{sec:setup} set context and methodology.
\S\ref{sec:characterization} tests which dimensions predict opportunity.
\S\ref{sec:routing} reports routing evaluation.
